{
\begin{fileviewerbox}{Configuration (.yaml)}
\begin{Verbatim}[breaklines=true]
# Prompt Templates: Control how observations of environment are shown to agent
system_template: |
  First `system` message shown to agent
instance_template: |-
  Instance prompt, contains task instance-specific content
next_step_template: |-
  Format template of per-turn observation (Contains standard output from agent's action)
next_step_no_output_template: |-
  Format template of observation when there is no standard output from the agent's action
format_error_template: |-
  Format template of error message (Used when agent's action causes an error)
demonstration_template: |
  Format template for showing a demonstration to the agent
demonstrations:
- `trajectories/<username>/<experiment folder>/*.traj` 
- File is a demonstration of how to solve a task. This could an agent generated trajectory.
- You can include 1+ demonstrations

# Environment States: Define features of the SWEEnv environment
env_variables:
# Default variables for SWEEnv at the beginning of each instance
  CURRENT_FILE: 0
  CURRENT_LINE:
  OVERLAP:
  SEARCH_FILES:
  SEARCH_INDEX:
  SEARCH_RESULTS:
  WINDOW_SIZE:
  START_INDEX:
  END_INDEX:
  START_CURSOR:
  END_CUROSR:
  START_CURSORS_MARK:
  END_CURSOR_MARK:
state_command: |
# `state_command` allows you to update state variables to reflect any aspect of the environment (e.g. current working directory)
  name: state
  code: |
    state() { echo '{"pwd": "'$PWD'"}';

# Action Interface: Define how an agent interacts with the SWEEnv environment
command_files:
- path/to/bash_file.sh
- Each file contains a list of commands implemented in bash
- You can include 1+ command files
parse_command: Reference to functionality for defining command documentation
history_processor: Reference to functionality for controlling agent's message history
parse_function: Parser run on agent output
\end{Verbatim}
\end{fileviewerbox}
\captionsetup{type=figure}
\captionof{figure}{
An example of the configuration file that defines the SWE-agent ACI.
A configuration is represented as a single \texttt{.yaml} file, allowing you to
define the commands that agents may use,
write prompts shown to the agent over the course of a single trajectory,
and control the input/output interface that sits between the agent and environment.
}
\label{fig.prompts.example_configuration}
}